\section{Production Deployment Architecture}
\label{sec:deployment}

This section addresses an obvious reviewer concern: the headline
results of \S\ref{sec:evaluation} are produced by a discrete-event
simulator (\S\ref{sec:simulator}), not by a deployed system. The
validity of conclusions like the Wait-Awhile +305\% inflation and
the do-no-harm property rests on the simulator faithfully capturing
production FaaS dynamics. This section specifies a concrete
production-validation architecture on Knative (the de facto
open-source FaaS platform), identifies the semantic mismatches with
our simulator, and proposes mitigations. We have not deployed this
architecture; the discussion is part deployment plan, part
limitations statement.

\subsection{Component Mapping: Simulator to Knative}
\label{sec:deployment:mapping}

\paragraph{Break-even gate (Lemma~\ref{lemma:tradeoff}).}
The closed-form break-even rate $\lambda^*$ that gates region
admission becomes a Kubernetes Custom Resource maintained by a
custom Operator. The Operator periodically (every $\sim 60$\,s) polls
ElectricityMaps via the same fetch protocol used in
\code{scripts/fetch\_electricitymaps.py}, queries Prometheus for the
arrival-rate estimator $\hat\lambda$ produced by the
\code{ArrivalRateTracker} of Algorithm~\ref{alg:greenfaas}, and
publishes a per-region \code{ConfigMap} indicating whether the
spatial or temporal regime is active. The Operator never makes
per-invocation decisions; it only updates platform state at
$\sim$minute granularity, leaving per-invocation routing to the
proxy below.

\emph{Caveat on poll latency.} A 60-second Operator poll interval
means up to one minute of staleness in the regime classification
during rapid carbon-intensity shifts (e.g., a wind front passing
through Germany at midnight). This produces up to one minute of
mis-routing for invocations dispatched between an intensity change
and the next Operator update. The bound is benign for slow-varying
grids (where intensity rarely moves more than 5\% per minute) but
worth instrumenting in production: log the Operator update timestamps
to OpenTelemetry alongside the routing decisions, so post-hoc
attribution of mis-routing to poll latency is possible.

\paragraph{Per-invocation routing.}
The natural injection point for per-invocation \GreenFaaS{}
decisions is \emph{not} the Kubernetes Scheduler (which routes pods
to nodes within a cluster, not invocations to clusters across
regions). Instead, we propose a sidecar in the Knative
\code{Activator} flow. The Activator already sits between the Knative
\code{Service} ingress and the function pods, mediating cold starts;
adding a \GreenFaaS{} decision pass before the Activator forwards
the request is mechanically straightforward.

The proxy reads the \code{ConfigMap}, computes the per-invocation
score from Algorithm~\ref{alg:greenfaas}, and either (a) holds the
request locally (temporal deferral, $\Delta t \le$ class-dependent
budget) and forwards to the local Activator when the deferral
completes, or (b) rewrites the request URL to the alternate cluster's
Activator and forwards. The local Activator then handles cold-start
detection and pod-warming as it does today; our policy only changes
\emph{which} cluster sees the request.

\paragraph{Cross-cluster forwarding.}
Cross-cluster forwarding requires that all participating clusters
share a function-image registry and that requests carry sufficient
identifying information (function name, request payload). Knative's
existing multi-tenant model handles this through the
\code{Service.serving.knative.dev/cluster-local} annotation; for
multi-region deployment, each region runs an identically-named
Knative \code{Service}, and the proxy targets the alternate region's
ingress URL. No changes to Knative itself are required.

\subsection{Replacing the Power Model}
\label{sec:deployment:power}

The simulator's $P_a = 4$\,W active, $P_i = 0.3$\,W idle is a
calibration to typical container-level draws on a server-class CPU.
Production validation requires \emph{measured} power, not assumed.

\paragraph{Kepler.}
\code{Kepler} (Kubernetes-based Efficient Power Level Exporter)
\citep{kepler} uses eBPF to attribute CPU and memory energy
consumption to individual pods on the Kubelet. A pod's energy draw
during execution is measured at sub-second resolution and published
to Prometheus. For warm-but-idle pods, Kepler measures the residual
draw (typically dominated by background runtime and language-VM
overhead rather than the function code itself), giving an empirical
$P_i$ that may differ from our 0.3\,W assumption by a substantial
factor depending on the runtime.

The relevance for our paper is that $\beta = P_i T_w / (P_a \tau_e)$
is the dimensionless quantity governing both the
Lemma~\ref{lemma:tradeoff} break-even and
Conjecture~\ref{conj:lech-lb}'s asymptotic lower bound;
production-measured $\beta$ via Kepler is what determines the realized
empirical behavior. If Kepler-measured $\beta$ on the chosen Knative
runtime differs from our 150 by a factor of 2--3, the qualitative
conclusions of \S\ref{sec:evaluation} should reproduce, but the
quantitative magnitudes will shift.

\paragraph{Kepler measurement caveats.} eBPF-based per-pod power
attribution has known accuracy limitations: published Kepler
benchmarks report 5--15\% per-pod attribution error depending on
workload type (compute-bound functions show lower error than
I/O-bound ones; co-resident pods sharing a CPU socket are harder to
disentangle). This uncertainty propagates directly to $\beta = P_i
T_w / (P_a \tau_e)$: a 10\% error on $P_i$ produces a 10\% error on
$\beta$, which in turn produces a similar-magnitude error on the
Lemma~\ref{lemma:tradeoff} break-even rate $\lambda^*$. A production
deployment should therefore report $\beta$ and $\lambda^*$ with
explicit confidence intervals derived from Kepler's measurement
uncertainty, and verify that the qualitative conclusions of
\S\ref{sec:evaluation} are robust to $\beta$ varying within
$[100, 200]$.

\paragraph{Carbon attribution.}
Per-pod energy is converted to carbon by multiplying by the
ElectricityMaps-reported $C(t)$ of the cluster's region at execution
time. This is the same model the simulator uses; the only
difference is that the energy term is measured rather than computed.

\subsection{Latency and SLA Instrumentation}
\label{sec:deployment:latency}

The simulator tracks latency as $\mathrm{Lat}_k = (s_k - t_k) +
\tau_e + \mathbf{1}[\text{cold}] \tau_c + \mathrm{RTT}(r_k)$. Production
latency includes everything the simulator omits: ingress overhead,
TLS handshake variance, pod scheduling jitter, image-pull delay if
the image is not pre-pulled, language-VM cold-start variance, and
network jitter on cross-region forwards. SLA violations could
plausibly come from any of these sources, not just the
\GreenFaaS{}-induced deferral that the simulator models.

\paragraph{OpenTelemetry traces.}
The proxy sidecar emits distributed traces tagged with
\GreenFaaS{}'s decision (commit time, target region, cold/warm
classification at decision time) and the realized latency components
above. The trace collection backend stores these per-invocation, and
the analysis pipeline computes the SLA violation rate, attributed by
component. This separates ``\GreenFaaS{} caused the violation by
deferring too long'' from ``\GreenFaaS{} forwarded to a region where
the Activator had a cold-start spike unrelated to our policy.''

\subsection{Semantic Mismatch: $T_w$ Granularity}
\label{sec:deployment:mismatch}

The most consequential semantic mismatch between our simulator
(and Lemma~\ref{lemma:tradeoff}) and Knative is the granularity of
the warm-pool TTL. The simulator models $T_w$ as a per-container
property: each warm container expires exactly $T_w$ seconds after
its last invocation. Knative's KPA (Knative Pod Autoscaler)
maintains $T_w$ as a per-\code{Revision} property
(\code{scale-to-zero-grace-period}), not per-pod, and the
scale-to-zero decision is made by a controller poll at the configured
interval, not exactly at the TTL boundary.

The practical consequence: Knative's effective $T_w$ for a
\emph{specific} pod can range from $T_w$ to $T_w + \delta$ where
$\delta$ is the controller poll interval (default 60\,s). The
$\lambda^*$ computed by Lemma~\ref{lemma:tradeoff} is therefore an
\emph{approximation} of the production-realized break-even rate; the
approximation error is bounded by $\delta / T_w$, around 10\% for
default Knative settings.

\paragraph{Mitigation.}
Two production configurations tighten this:
(i) set Knative's \code{container-concurrency: 1} so each warm pod
serves one invocation at a time (matching our single-container
model exactly), and (ii) reduce the controller poll interval to
$\le 5$\,s. Together these bring the simulator-vs-production
$T_w$ mismatch below 1\%.

\emph{Caveat on aggressive polling at scale.} A 5-second poll
interval increases API-server load proportionally; on clusters with
$>$1000 \code{Revision} objects, the load increase can require
horizontal scaling of the \code{kube-apiserver} and add measurable
latency to other control-plane operations. A production deployment
at scale should benchmark the API-server load impact and consider
intermediate values (e.g., 15\,s) that retain most of the $T_w$
fidelity while keeping control-plane overhead acceptable. For
\GreenFaaS{}'s purposes, even a 30-second poll interval keeps the
$T_w$ mismatch under 5\%, which is small compared to the
Kepler-induced uncertainty on $\beta$ noted in
\S\ref{sec:deployment:power}.

\subsection{Cold-Start Variance}
\label{sec:deployment:coldstart}

The simulator treats $\tau_c$ as a constant per-function property
(typically $1$\,s in our catalog). Knative cold starts have orders
of magnitude more variance, driven by: image-pull time from the
container registry (10\,ms--30\,s depending on cache state),
runtime initialization (sub-second for compiled languages,
multi-second for Python/Node with heavy dependency graphs), and
network setup (typically 100--500\,ms).

\paragraph{Mitigation.}
A \code{DaemonSet} that pre-pulls function images to every node
removes the dominant variance source. With pre-pulled images and a
fast runtime (Go or compiled Node), production $\tau_c$ converges to
$200$--$500$\,ms with low variance, close to our simulator's
$1$\,s default. The simulator's $\tau_c$ should be re-calibrated
against Kepler-measured cold-start durations once the deployment is
running; if production $\tau_c$ differs significantly from $1$\,s,
the $\alpha = \tau_c / \tau_e$ ratio in Lemma~\ref{lemma:tradeoff}
shifts and $\lambda^*$ should be recomputed.

\subsection{What This Validates and What It Does Not}
\label{sec:deployment:validates}

A production deployment built as above would validate:
\begin{itemize}
\item Whether the Wait-Awhile and Lechowicz failure modes
(\S\ref{sec:evaluation:lechowicz}) reproduce on a real platform
with real Kepler-measured energy.
\item Whether the do-no-harm property
(\S\ref{sec:evaluation:fine-donoharm}) holds when production
$T_w$ has the controller-poll jitter described above.
\item Whether the \GreenFaaS{}-Spatial gap stays within our
empirical 0.058pp $\pm$ 0.032pp range
(\S\ref{sec:evaluation:realboth}) under production conditions.
\end{itemize}

It would \emph{not} resolve Conjecture~\ref{conj:lech-lb}, which
requires a formal multi-container model and an adversary
construction; a production deployment can provide empirical support
but not a proof.

The deployment also exposes a class of issues the simulator does not
model: cross-region network partitions, controller failures during a
scale-to-zero event, and registry availability for cold-start
image-pulls. These are infrastructure-reliability issues rather than
scheduling-policy issues; we treat them as orthogonal.

\paragraph{Why we have not done this.}
A full production validation requires multi-region Kubernetes
clusters with Knative installed, Kepler integration, an
ElectricityMaps subscription with sufficient quota, and several
days of running representative workloads. Our build environment does
not provide cluster resources, and the validation is a non-trivial
engineering project rather than a paper-revision task. We treat the
architecture described above as a concrete plan for follow-up work,
and note it as the most consequential validation step still
remaining.
